% Generated draft: requires mathematical and source review.
\documentclass[11pt]{article}
\usepackage[margin=25mm]{geometry}
\usepackage{fontspec}
\setmainfont{Latin Modern Roman}
\usepackage{amsmath,amssymb,mathtools,unicode-math}
\setmathfont{Latin Modern Math}
\usepackage{enumitem,needspace,xurl,hyperref}
\setlength{\parindent}{0pt}
\setlength{\parskip}{5pt}
\setlength{\emergencystretch}{3em}
\newcommand{\N}{\mathbb{N}}
\newcommand{\Z}{\mathbb{Z}}
\newcommand{\Q}{\mathbb{Q}}
\newcommand{\R}{\mathbb{R}}
\newcommand{\C}{\mathbb{C}}
\newcommand{\F}{\mathbb{F}}
\newcommand{\A}{\mathbb{A}}
\newcommand{\PP}{\mathbb{P}}
\newcommand{\E}{\mathbb{E}}
\newcommand{\eps}{\varepsilon}
\newcommand{\dd}{\,\mathrm{d}}
\begin{document}
\section*{UnsolvedMath ID: 1100113}
\textbf{Source problem code:} catalogue entry
\subsection{Definitions and mathematical statement}
Swarup's source question is: prove a combination theorem for relatively hyperbolic groups. Its intended general shape can be recorded as follows. Let $\mathcal G$ be a graph of groups over a connected graph $\Lambda$, with vertex groups $G_v$, edge groups $G_e$, and injective edge-to-vertex homomorphisms. Equip each vertex group with a peripheral family $\mathcal P_v$ such that $(G_v,\mathcal P_v)$ is relatively hyperbolic in the chosen sense. Seek suitable hypotheses on the edge groups, their embeddings, and the peripheral families, and an explicitly described peripheral family $\mathcal P$ in $G=\pi_1(\mathcal G)$, under which
\[
\bigl[(G_v,\mathcal P_v)\text{ is relatively hyperbolic for every vertex }v\bigr]
\ \text{together with the specified combination hypotheses}\ 
\Longrightarrow\ (G,\mathcal P)\text{ is relatively hyperbolic}.
\]
A one-edge graph yields either an amalgamated product or an HNN extension. This is a theorem schema, not an assertion that arbitrary amalgams of relatively hyperbolic groups satisfy the conclusion. The supplied one-sentence question fixes neither a definition of relative hyperbolicity nor the necessary edge hypotheses, finiteness assumptions, or resulting peripherals, so it does not determine one uniquely quantified combination theorem.

The supplied later summary labels the question answered affirmatively in several frameworks. It describes Dahmani (2003) as combining convergence-group actions along parabolic-type or ``liminal'' subgroups subject to geometric-finiteness and intersection-control conditions, with an assembled compactum and expected peripherals; it also reports the application to relative hyperbolicity of limit groups with respect to maximal noncyclic abelian subgroups. It describes Alibegovi\'c (2005) as treating a one-edge graph with a liminal edge group, an almost-malnormality condition, and an ``isolated''-type compatibility condition, with resulting peripherals from vertex peripherals not contained in the edge group. The exact conditions and identifications are not supplied and are not replaced here by guessed sufficient conditions.

The summary further identifies Mj--Reeves (2008) as a geometric combination theorem explicitly answering Swarup, but its description ends at ``a tree of strongly relatively'' and supplies no complete hypothesis or conclusion. A precise statement of that version cannot be recovered from the truncated record. The supplied solved label records these attributions rather than a theorem proved or independently verified in this draft.
\subsection{Short English statement}
The question asks when gluing relatively hyperbolic groups along subgroups again gives a relatively hyperbolic group, with an appropriate peripheral family. The supplied later summary reports answers by Dahmani, Alibegovi\'c, and Mj--Reeves in different frameworks. Their exact hypotheses are absent, and the last description is truncated, so the draft preserves a combination-theorem schema and the reported resolution without inventing a universal gluing theorem.
\subsection{Sources}
\begin{itemize}
\item Local source registry: \url{data/unsolvedmath/problems.json}, permanent UnsolvedMath ID $1100113$, source code catalogue entry, including the supplied research summary and solved label. Dataset attribution: Ulam AI and UnsolvedMath Contributors; source attribution and rights are retained.
\item Mladen Bestvina, \emph{Questions in Geometric Group Theory}, source-list date 2004, Question 1.13, attributed there to Swarup; source metadata identifies PDF page 3: \url{https://www.math.utah.edu/~bestvina/eprints/questions-updated.pdf}. The record reports manual transcription and access on July 29, 2026, with current-status and rights review flags.
\item The supplied research summary names Dahmani (2003), Alibegovi\'c (2005), and Mj--Reeves (2008) for distinct combination theorems. It supplies no full titles, precise theorem numbers, or separate bibliographic links for those references. The linked source and these literature claims were not independently inspected.
\end{itemize}
\end{document}
